%قبل از شروع مسئله ترجیحاً چیزی ننویسید. اما اگر خواستید بنویسید، در قالب لیستی (فقرات و شمارش) بنویسید.
%برای اضافه کردن نمره، بعدی از شروع‌مسئله، [x] را اضافه کنید و به جای x  نمره مورد نظر را قرار دهید. برای زیرمسئله‌ها هم به همین صورت است.
%\شروع‌مسئله[10]
\شروع‌مسئله[26]
% صورت اصلی مسئله را در این‌جا بنویسید.
در میان الگوریتم‌هایی که در اختیارتان قرار گرفته است، اثباتی برای سریع‌ترین الگوریتم درست ارائه دهید. در ادامه، یک اثبات مناسب برای درستی الگوریتم چهارم آورده شده است. شما می‌توانید از الگوی این اثبات برای اثبات خود استفاده کنید.

\بدون‌تورفتگی
\متن‌سیاه{اثبات الگوریتم چهارم.}
اثبات درستی الگوریتم ۴ را در دو گام ارائه می‌دهیم. در گام اول ثابت می‌کنیم که الگوریتم ۴ همواره یک خانه را برمی‌گرداند و سپس در گام دوم اثبات می‌کنیم اگر خانه‌ای برگردانده شود حتما قله است. به‌وضوح این دو گزاره درستی الگوریتم را نتیجه می‌دهند.
\شروع{شمارش}[1. ]
	\فقره \textbf{این الگوریتم حتما یک خانه را برمی‌گرداند.}
	برای اثبات درستی این گزاره فرض کنید جدول موجود در یک مرحله $m\times n$ است. حال در هر مرحله بسته به این که تقسیم را از سطر انجام می‌دهیم یا از ستون، ابعاد مسئله به یکی از چهار حالت 
	$ \lfloor \frac{m}{2}\rfloor \times n$،
	$  (m - \lfloor \frac{m}{2}\rfloor - 1) \times n$،
	$ m \times \lfloor \frac{n}{2}\rfloor$ و یا 
	$ m \times (n - \lfloor \frac{n}{2}\rfloor - 1)$  
	تبدیل می‌شود. توجه کنید که مقدار $m$ یا $n$ در هر مرحله کم می‌شود و یا این‌که مقداری توسط الگوریتم به عنوان خروجی برگردانده می‌شود. ادعا می‌کنیم که $m$ و $n$ همواره مثبت باقی می‌مانند. بنابر کاهش $m$ یا $n$  در هر مرحله، اگر این ادعا را ثابت کنیم، نشان داده‌ایم که الگوریتم حتماً به اتمام می‌رسد.
	
	برای اثبات این ادعا از برهان خلف استفاده می‌کنیم و فرض می‌کنیم که $n \leq 0$  (استدلال مشابه برای $m$ برقرار است). ابتدا توجه کنیم که اگر $n > 2$   باشد، مقدار $n$ در مرحلۀ بعد نیز مثبت می‌ماند.  پس کافی است این موضوع را تنها برای
	 $n = 1$
	 و 
	 $n =2$
	بررسی کنیم. همچنین می‌توانیم فرض می‌کنیم در حالت تقسیم ستونی هستیم زیرا در غیر این صورت مقدار $n$ تغییر نمی‌کند. 
\renewcommand{\labelitemi}{{\scriptsize\textbullet{}} }
	\شروع{فقرات}
	\فقره
	اگر $n=1$  باشد، آن‌گاه محاسبۀ بیشینۀ ستون مرکزی در واقع محاسبۀ بیشینۀ کل جدول است که بنابر تعریف یک قله است و برنامه آن را باز می‌گرداند (خانۀ \lr{best\_seen} نیز در همین ستون است، پس مقدار در این خانه بیش‌تر مساوی مقدار \lr{best\_seen} نیز هست). 
	\فقره
	در حالت $n=2$ نیز  دو حالت ممکن است. اول اینکه خانۀ بیشینۀ ستون وسط، یک قله باشد و الگوریتم آن را برگرداند. حالت دوم این است که الگوریتم را برای مسئلۀ کوچک‌تر شامل \lr{best\_seen} صدا کنیم (اگر مقداری باز نگردد، \lr{best\_seen} حتماً در ستون دیگر قرار دارد). در این حالت نیز $n=1$ می‌شود.
	\پایان{فقرات}
	 پس در دو هر دو حالت $n$  مثبت می‌ماند، که نشان می‌دهد فرض خلف باطل است.
	\فقره
	\متن‌سیاه{ اگر خانه‌ای برگردانده شود حتماً یک قله (در مسئلۀ اصلی) است.} فرض می‌کنیم الگوریتم خانۀ $(r_1, c_1)$ را بازگردانده باشد. توجه کنید این خانه در زیرمسئلۀ خودش حتما یک قله است. حال از برهان خلف استفاده می‌کنیم و فرض می‌کنیم این خانه در مسئلۀ اصلی قله نباشد.  دنبالۀ مسائل بازگشتی را در نظر می‌گیریم. کوچکترین مسئله را در این دنباله درنظر بگیرید که خانۀ $(r_1, c_1)$ در آن یک قله نیست. در این مرحله، این خانه همسایه‌ای مانند $(r_2, c_2)$ دارد که مقدار آن بزرگ‌تر از مقدار پیدا شده است، یعنی 
	$val(r_2, c_2) > val(r_1, c_1)$. همچنین خانۀ  
	$(r_2, c_2)$
	در زیرمسئلۀ بعدی وجود ندارد؛ پس این خانه حتما روی سطر یا ستونی بوده است که در آن مرحله از روی آن تقسیم انجام شده است. امّا الگوریتم، مقدار \lr{best\_seen} را با توجه به ستون‌ها و سطر‌های تقسیم‌شونده مقداردهی می‌کند. اگر خانۀ \lr{best\_seen} را $(r_{bs}, c_{bs})$ بگیریم داریم 
	$val(r_2, c_2) \leq val(r_{bs}, c_{bs})$.
	همچنین از آن‌جایی که الگوریتم مقدار  
	$val(r_1, c_1)$
	را برگردانده است، می‌دانیم 
		$val(r_{bs}, c_{bs}) \leq val(r_{1}, c_{1})$.
		پس داریم
		\begin{equation*}
			val(r_{bs}, c_{bs}) \leq val(r_1, c_1) < val(r_2,c_2) \leq val(r_{bs}, c_{bs})
		\end{equation*}
		که تناقض است. پس فرض خلف باطل است و حکم ثابت می‌شود. \qed
\پایان{شمارش}

%اگر مسئله چند قسمت داشت به این صورت عمل کنید:
%\شروع{زیرمسائل}
%\زیرمسئله[0] % مسئلۀ ۱آ
%\زیرمسئله % مسئلۀ ۱ب
%\زیرمسئله % مسئلۀ ۱پ
%\زیرمسئله % مسئلۀ ۱ت
%\پایان{زیرمسائل}

\پاسخ{
	برای اثبات  مسئله از الگوی صورت سؤال استفاده می‌کنیم. برای راحتی ابتدا یک تعریف انجام می‌دهیم. قاب را خانه‌های روی سطر‌های اول، وسط و آخر و همچنین ستون‌های اول، وسط و آخر در نظر می‌گیریم. حال دو گزارۀ بالا را برای این الگوریتم ثابت می‌کنیم.
	\شروع{شمارش}
	\فقره \textbf{این الگوریتم حتماً یک خانه را برمی‌گرداند.}
	برای اثبات درستی این گزاره فرض می‌کنیم اندازۀ مسئلۀ اولیه $n\times n$ است. می‌توان دید در هر مرحله ابعاد مسئله به یکی از دو حالت زیر تبدیل می‌شود:
	$ (\lfloor \frac{n + 1}{2} \rfloor - 2) \times (\lfloor \frac{n + 1}{2} \rfloor - 2)$ و یا 
	$ (\lceil \frac{n + 1}{2} \rceil  - 2)\times (\lceil \frac{n + 1}{2} \rceil - 2)$.
	 توجه کنید که مقدار ابعاد ماتریس در هر مرحله کم می‌شود و یا این‌که مقداری توسط الگوریتم برگردانده می‌شود. ادعا می‌کنیم که $n$ همواره مثبت باقی می‌ماند. بنا بر توضیحات قبلی اگر این ادعا را ثابت کنیم، نشان داده‌ایم که الگوریتم حتماً یک مقدار را برمی‌گرداند.  
	 
	توجه کنیم که اگر $n > 4$ باشد، $n$ حتما در مرحلۀ بعد نیز مثبت باقی خواهد ماند. همچنین اگر $n = 4$ باشد، یا در قاب یک قله پیدا می‌کنیم یا $n=1$ خواهد شد (دقت کنیم در هر مرحله، زیرماتریسی که شامل همسایۀ بزرگ‌تر است، انتخاب می‌شود، همچنین این همسایۀ برزگ‌تر در قاب نیست). اگر $n \leq 3$ باشد، آن‌گاه تمامی خانه‌های جدول در قاب قرار می‌گیرند و بیشینۀ قاب، بیشینۀ جدول و قله خواهد بود که بازگردانده می‌شود. پس همواره مقدار $n$ مثبت باقی می‌ماند که نشان می‌دهد حتماً یک خانه توسط این الگوریتم بازگردانده می‌شود.
	\فقره
	\textbf{ اگر خانه‌ای برگردانده شود حتما ًیک قله (در مسئلۀ اصلی) است.} فرض می‌کنیم الگوریتم خانۀ $(r_1, c_1)$ را بازگردانده باشد. توجه کنید این خانه در زیرمسئلۀ خودش حتما ًیک قله است و همچنین این خانه بیشینۀ قاب زیرمسئلۀ خود است. حال اگر این خانه بر روی محیط قاب نباشد (چهار ضلع قاب) که قله بودن در این زیرمسئله و مسئلۀ اولیه معادل است و این خانه در مسئلۀ اصلی هم قله است. پس فرض می‌کنیم که این خانه در محیط قاب باشد. 
	از برهان خلف استفاده می‌کنیم و فرض می‌کنیم همسایه‌ای مانند 
$(r_2, c_2)$	
در مسئلۀ اصلی باشد که مقدار آن از مقدار 
$(r_1, c_1)$
برزگ‌تر باشد.  توجه کنید این خانه لزوماً در قاب مسئلۀ پیشین قرار دارد. بیشینۀ قاب مسئلۀ پیشین را 
$g = val(r_3, c_3)$
  می‌گیریم. همچنین می‌دانیم این خانه همسایه‌ای داشته که مقدارش از آن بیشتر بوده و حتماً در قاب مسئلۀ فعلی قرار دارد؛ این خانه را   $(r_4, c_4)$
  می‌نامیم. می‌دانیم 
  مقدار
  $(r_1, c_1)$
  بیشینه در این قاب بوده است، یعنی 
  $val(r_1, c_1) \geq val(r_4, c_4)$.
پس داریم
\begin{equation*}
	val(r_1, c_1) \geq val(r_4, c_4) > val(r_3, c_3) = g \geq val(r_2, c_2)
\end{equation*}
که نشان می‌دهد فرض خلف‌مان غلط است. پس حکم ثابت می‌شود. \qed
\پایان{شمارش}
}





%%%%%%%%%%%%%%%%%%%%%%%%%%%%%%%%%%%%%%%%%%%%%%%%%%%%%
% برخی دستورات مفید برای تمرین‌ها. %
%%%%%%%%%%%%%%%%%%%%%%%%%%%%%%%%%%%%%%%%%%%%%%%%%%%%%

% فرمول‌های نمونه:
%$f(x) = \Theta(x)$
%$T(x, y) \leq \log(x) + 2^y + \binom{2n}{n}$
% {\tt code\_function}


% لیست موارد بدون شماره‌گذاری:
%\شروع{فقرات}
%    \فقره مورد اول 
%        \شروع{فقرات}
%            \فقره مورد اول داخلی 
%                \شروع{فقرات}
%                    \فقره اولین مورد خیلی داخلی 
%                    \فقره دومین مورد خیلی داخلی 
%                \پایان{فقرات}
%            \فقره مورد دوم داخلی 
%        \پایان{فقرات}
%    \فقره مورد دوم 
%\پایان{فقرات}

% لیست موارد شماره‌گذاری شده:
%\شروع{شمارش}
%    \فقره یک مورد 
%    \فقره یک مورد دیگر 
%    \فقره باز هم یک مورد دیگر
%\پایان{شمارش}

% برای نوشتن یک فرمول چند خطی که یک شماره بخوره. برای حالت‌های دیگر به سایت زیر مراجعه کنید:
% https://www.overleaf.com/learn/latex/Aligning_equations_with_amsmath
%\begin{align} 
%    \begin{split}
%        (x+y)^3 &= (x+y)^2(x+y) \\
%                &= (x^2+2xy+y^2)(x+y) \\
%                &= (x^3+2x^2y+xy^2) + (x^2y+2xy^2+y^3) \\
%                &= x^3+3x^2y+3xy^2+y^3
%    \end{split}                                 
%\end{align}

% نوشتن ماتریس و آرایه دوبعدی
%\begin{align}
%    A = 
%    \begin{bmatrix}
%        A_{11} & A_{21} \\
%        A_{21} & A_{22}
%    \end{bmatrix}
%\end{align}

% اضافه کردن تصویر با شرح. برای اضافه کردن تصویر بدون شرح، فقط دستور درج‌تصویر را استفاده کنید.
%\شروع{شکل}[h!]
%\تنظیم‌ازوسط
%  \درج‌تصویر[پهنا=0.5\پهنای‌متن]{img.jpg}
%\شرح{شرح تصویر}
%\برچسب{تصویر:۱}
%\پایان{شکل}
%\پایان{زیرمسائل}

%کد زیر را ببینید.
%%\lstinputlisting{source_filename.py}
%\begin{latin}
%\begin{lstlisting}
%# Import the modules
%import sys
%import random
%
%ans = True
%
%while ans:
%question = raw_input("Ask the magic 8 ball a question: (press enter to quit) ")
%
%answers = random.randint(1,8)
%
%if question == "":
%sys.exit()
%
%elif answers == 1:
%print "It is certain"
%
%elif answers == 2:
%print "Outlook good"
%
%elif answers == 3:
%print "You may rely on it"
%
%elif answers == 4:
%print "Ask again later"
%
%elif answers == 5:
%print "Concentrate and ask again"
%
%elif answers == 6:
%print "Reply hazy, try again"
%
%elif answers == 7:
%print "My reply is no"
%
%elif answers == 8:
%print "My sources say no"
%\end{lstlisting}
%\end{latin}
